\section{Conclusão}


\begin{frame}
	\frametitle{Principais contribuições}	
	\begin{itemize}
		\item Construção do dataset: integração, organização e curadoria de diferentes bases públicas.
		\item Avaliação de modelos: comparação de arquiteturas e análise da generalização.
		\item Explicabilidade: utilização do Grad-CAM para investigar as regiões consideradas pelos modelos.
		\item Aplicação prática: disponibilização do modelo em uma aplicação Android.
	\end{itemize}

\end{frame}


\begin{frame}
	
	\frametitle{Conclusão e trabalhos futuros}	
	\begin{itemize}
		
		\item Conclusões
		\begin{itemize}
			\item A integração e curadoria das bases foram importantes para construir um conjunto mais consistente.
			\item Os experimentos permitiram comparar diferentes arquiteturas e avaliar suas limitações.
			\item A análise com Grad-CAM ajudou a compreender o comportamento dos modelos.
			\item A aplicação Android demonstrou a possibilidade de utilizar o modelo fora do ambiente de treinamento.
		\end{itemize}
		
		\item Possíveis trabalhos futuros:
		\begin{itemize}
			\item Ampliar a diversidade de imagens.
			\item Melhorar a generalização para diferentes ambientes.
			\item Investigar estratégias para reduzir erros de classificação.
			\item Aprimorar a aplicação Android.		 	
		\end{itemize}
		
	\end{itemize}
	
	
\end{frame}
